% ------------------------------------------------------------------ 6
\section{Exact paradigm 2: branch-and-bound}\label{sec:bnb}

The DP is exponential in \(k\) and useless beyond \(k\approx20\). Branch-and-bound is exponential in the number of branching
decisions instead, so it helps when a good lower bound settles most decisions before any branching happens. Our design follows the
structure of the strongest exact Steiner solvers (Koch and Martin 1998; Polzin and Vahdati Daneshmand 2001), cut down to what can be
written without an LP solver: exact graph reductions, a dual ascent lower bound, a heuristic upper bound, reduced-cost
elimination, and best-first branching on Steiner vertices.

\subsection{Search space and branching}
A search node fixes some non-terminals as \emph{included} (they become terminals) and some vertices and edges as \emph{excluded}
(they are deleted). Branching on a free non-terminal \(v\) creates two children, \(v\) included and \(v\) excluded. Every Steiner tree
either contains \(v\) or does not, so the two children together keep every feasible tree. Each branching fixes one vertex, so every
root-to-leaf path has at most \(n-k\) branchings, and the search tree is finite with at most \(2^{n-k+1}\) nodes. We branch on the free
non-terminal with the smallest \(d'(r,v)+d'(v,T)\) (defined in \cref{sec:elim}): the vertex that the dual solution considers cheapest
to use, which is most likely in an optimal tree and therefore tightens the bound quickly in the excluding child.

Open nodes wait in a Fibonacci heap keyed by their lower bound, and the node with the smallest bound is expanded next
(best-first search). The minimum key is a global lower bound on \(\OPT\): as soon as it reaches the incumbent cost \(U\) the search
stops with a proof of optimality.

\subsection{Exact reductions}
Before the search the graph is reduced. All four tests keep at least one optimal tree.

\begin{lemma}\label{lem:reductions}
Let \(k\ge2\). (R1) A non-terminal of degree at most 1 lies in no optimal tree with only terminal leaves; delete it.
(R2) If a terminal \(t\) has a single edge \(\{t,u\}\), that edge is in every feasible tree; add \(w(t,u)\) to a fixed cost, delete \(t\)
and make \(u\) a terminal. (R3) If a non-terminal \(v\) has exactly two edges \(\{a,v\},\{v,b\}\), replace them by one edge
\(\{a,b\}\) of weight \(w(a,v)+w(v,b)\), keeping the cheaper one if \(\{a,b\}\) already exists. (R4) Every edge with
\(w(u,v)>d(u,v)\) can be deleted, all at once.
\end{lemma}
\begin{proof}
(R1) A non-terminal leaf can be pruned without raising the cost. (R2) \(t\) must be connected to another terminal, and its only
edge is the only way. (R3) In a tree with only terminal leaves, \(v\) either does not appear or has degree 2 and uses both edges, which
is exactly the new edge; if a cheaper parallel edge exists, any tree using the path can switch to it. (R4) If a tree uses
\(\{u,v\}\), replace that edge by a shortest \(u\)--\(v\) path. The result is connected, spans \(T\) and is strictly cheaper, so the
edge is in no optimal tree. Edges that are strictly longer than the distance between their ends lie on no shortest path at all, so
deleting them changes no distance and the argument holds for all of them together.
\end{proof}
\noindent R1--R3 run until nothing changes, then R4 (one Dijkstra per vertex), and the whole cycle repeats while anything changes.
Each reduced edge remembers the list of original edges it stands for, so the final tree is mapped back to the input graph and checked
there.

\subsection{Lower bound: dual ascent on the cut formulation}
Orient the problem from the root: replace each edge by two opposite arcs of the same weight. A Steiner tree oriented away from \(r\) is
an arborescence \(A\), and \(x=\chi^{A}\) satisfies
\[\min\ \sum_{a}c_a x_a\quad\text{s.t.}\quad x(\delta^-(W))\ge1\ \ \text{for all }W\subseteq V,\ r\notin W,\ W\cap T\neq\emptyset,\qquad x\ge0, \tag{P}\]
where \(\delta^-(W)\) is the set of arcs entering \(W\). The dual is
\[\max\ \sum_W y_W\quad\text{s.t.}\quad \sum_{W:\,a\in\delta^-(W)}y_W\le c_a\ \ \text{for all arcs }a,\qquad y\ge0. \tag{D}\]

\begin{lemma}[admissibility]\label{lem:lb}
For a dual-feasible \(y\) with reduced costs \(c'_a=c_a-\sum_{W:a\in\delta^-(W)}y_W\ge0\), every Steiner arborescence \(A\) satisfies
\(c(A)\ge\sum_W y_W+c'(A)\). In particular \(\mathrm{LB}=\sum_Wy_W\le\OPT\).
\end{lemma}
\begin{proof}
\(c(A)=\sum_{a\in A}c'_a+\sum_{a\in A}\sum_{W:a\in\delta^-(W)}y_W=c'(A)+\sum_W y_W\,|A\cap\delta^-(W)|\), and
\(|A\cap\delta^-(W)|\ge1\) because \(A\) contains a path from \(r\notin W\) to a terminal in \(W\).
\end{proof}

Wong's dual ascent builds such a \(y\) greedily. For an active terminal \(t\), let \(W(t)\) be the set of vertices that reach \(t\)
through arcs of reduced cost 0. While \(r\notin W(t)\), raise \(y_{W(t)}\) by the smallest reduced cost on an arc entering \(W(t)\);
that arc becomes saturated (\(c'=0\)), its tail joins \(W(t)\), and the bound grows by the same amount. Reduced costs never go
negative, so \(y\) stays feasible and \cref{lem:lb} applies at every moment. Our implementation grows one component at a time and keeps
the arcs entering \(W(t)\) in a Fibonacci heap. Raising \(y\) by \(\delta\) lowers every entering arc by \(\delta\), so instead of
touching them it adds \(\delta\) to an offset \(Y\) and keys each arc by its reduced cost plus the value of \(Y\) when it entered the heap:
the heap minimum is always the next arc to saturate. A component stops growing when it reaches \(r\) or another active terminal
(that terminal will connect both). Each growth step saturates one arc and arcs never unsaturate, so a call makes at most \(2m\) steps
per terminal, \(O(km\log m)\) in total. At the root node the ascent is repeated from up to 16 different root terminals and the best bound
is kept.

\subsection{Upper bound and reduced-cost elimination}\label{sec:elim}
At every node the shortest-path heuristic (\cref{sec:greedy}, G2) runs on the node graph and on the subgraph of arcs with reduced cost
0. That subgraph connects every terminal to \(r\), and trees inside it tend to be cheap because the dual has already paid for those
arcs. At the root it runs from up to 64 start terminals. The best tree is then improved by local search on its vertex set (try deleting
each Steiner vertex and adding each neighbouring vertex, recompute the MST, keep improvements). This search runs at the root, on every
node whose heuristic beats the incumbent, and on every 64th node.

\begin{lemma}[reduced-cost elimination]\label{lem:elim}
Let \(d'(x,y)\) be the shortest-path distance under the reduced costs \(c'\), \(d'(v,T)=\min_{t\in T}d'(v,t)\), and \(U\) the cost of
the incumbent. A non-terminal \(v\) with \(\mathrm{LB}+d'(r,v)+d'(v,T)\ge U\) lies in no tree cheaper than \(U\). An arc
\((u,v)\) with \(\mathrm{LB}+d'(r,u)+c'_{uv}+d'(v,T)\ge U\) lies in no arborescence cheaper than \(U\).
\end{lemma}
\begin{proof}
Let \(A\) be an arborescence with only terminal leaves that contains \(v\). It contains a path from \(r\) to \(v\) and, below \(v\), a path
from \(v\) to a terminal; the two share no arc. Since \(c'\ge0\), \(c'(A)\ge d'(r,v)+d'(v,T)\), and \cref{lem:lb} gives
\(c(A)\ge\mathrm{LB}+d'(r,v)+d'(v,T)\ge U\). The arc case is the same with the path \(r\to u\to v\to T\). An edge is deleted when
both of its arcs fail the test.
\end{proof}
\noindent Both \(d'\) are one Dijkstra run on the heap (from \(r\), and backwards from all terminals). The deletions are valid for the
whole subtree of the node, because \(U\) can only decrease, so children inherit them.

\begin{theorem}\label{thm:bnb}
If the search ends before the limits, the incumbent is optimal.
\end{theorem}
\begin{proof}
Invariant: either the incumbent is optimal, or some optimal tree is feasible for an open node. It holds at the start, and it is kept by
branching (the children partition the trees of the parent), by the reductions (\cref{lem:reductions}) and by the eliminations
(\cref{lem:elim}), which only remove trees that cost at least \(U\). A node is discarded only when \(\mathrm{LB}\ge U\), and by
\cref{lem:lb} it then holds no tree cheaper than \(U\). When the heap is empty, or its smallest key is at least \(U\), no open node can
hold a cheaper tree, so the incumbent is optimal. Weights are integers, so \(\mathrm{LB}\) is an integer and the comparison is exact.
\end{proof}

\paragraph{Complexity.} One node costs \(O(m)\) to build, \(O(km\log m)\) for the dual ascent, \(O(k(m+n\log n))\) for each heuristic run,
\(O(m+n\log n)\) for the two reduced-cost Dijkstra runs and \(O(n(m+n\log n))\) per local search pass when it runs. The number of nodes is
at most \(2^{n-k+1}\) in the worst case, so the worst-case running time is exponential, as it must be unless P\,=\,NP. The
measurements below show how far below that bound the search stays.

\subsection{Results}
\begin{table}[H]\centering\small\setlength{\tabcolsep}{4pt}
\begin{tabular}{@{}lrrrrrrrrr@{}}\toprule
Set & \# & solved & \(|V|\) left & \(|E|\) left & root gap \% & nodes: median & max & mean s & max s\\\midrule
B & 18 & 18 & 60\% & 67\% & 0.05 & 1 & 3 & 0.01 & 0.05\\
C & 20 & 19 & 79\% & 65\% & 0.61 & 3 & 23{,}088 & 17.17 & 300.02\\
P6E & 15 & 15 & 97\% & 98\% & 0.22 & 1 & 13{,}707 & 1.10 & 13.06\\
I080 & 100 & 100 & 86\% & 91\% & 2.31 & 10 & 835 & 0.07 & 0.53\\
Grid & 27 & 27 & 90\% & 89\% & 0.06 & 1 & 8 & 0.02 & 0.04\\
Random & 6 & 6 & 75\% & 77\% & 0.42 & 28 & 64 & 1.95 & 8.73\\
Adversarial & 13 & 13 & 93\% & 93\% & 0.00 & 1 & 1 & 0.00 & 0.00\\
\bottomrule\end{tabular}
\caption{Branch-and-bound with a 300\,s limit. \(|V|\), \(|E|\) left: share of the graph that survives the reductions. Root gap: \((\OPT-\mathrm{LB}_{\text{root}})/\OPT\) for the dual ascent bound at the root.}\label{tab:bnb}
\end{table}
\noindent Instances not solved within the limit, with the bounds proved when the search stopped:
\begin{center}\small\begin{tabular}{@{}lrrrrr@{}}\toprule
Instance & \(k\) & lower bound & best tree & OPT & nodes\\\midrule
c18 & 83 & 110 & 117 & 113 & 23{,}088\\
\bottomrule\end{tabular}\end{center}
\begin{figure}[H]\centering
\begin{tikzpicture}
\begin{loglogaxis}[width=0.42\linewidth,height=5.6cm,xlabel={DP time (s)},ylabel={B\&B time (s)},
  colorbar,colormap={kk}{color=(cG3) color=(cG1)},point meta=explicit,colorbar style={ylabel={\(k\)},font=\footnotesize,width=2.5mm},
  tick label style={font=\footnotesize},label style={font=\small},grid=major,grid style={black!8}]
\addplot[scatter,only marks,mark=*,mark size=1.5pt,fill opacity=0.7,draw opacity=0] table[x=dp,y=bb,meta=k,col sep=comma]{generated/dp_vs_bb.csv};
\addplot[black!50,dashed,domain=1e-5:200] {x};
\end{loglogaxis}
\end{tikzpicture}
\caption{The two exact methods on the 173 instances both solve. Below the diagonal B\&B is faster; it is faster on 97 of them.
Colour is the number of terminals: the DP time is set by \(k\), the B\&B time by how much of the graph the bound settles.}\label{fig:dpbb}
\end{figure}

